\section{Experiments}\label{sec:exp}

The experiments are organized around the three independently packaged components of the proposed system: link interval envelopes, the Lifelong Envelope Cache Tree (LECT), and the SafeBoxForest (SBF) query layer. This structure separates three questions that are otherwise easy to conflate. First, how expensive and how tight are the interval envelopes that turn a joint-space box into conservative workspace evidence? Second, how much of that evidence can be reused across scenes for a fixed robot? Third, when the evidence is consumed by a box forest, what build/query/path-quality trade-offs are obtained under strict path audit?

All planner-level success rates in this section use strict collision audit as the safety criterion. In particular, a Graph-of-Convex-Sets (GCS) solve is not counted as a successful plan unless the returned configuration-space path passes the same segment-wise audit used by SBF. This convention is important because provisional SBF boxes are conservative evidence carriers, not certified convex free-space regions. Directly treating every provisional box interior as a GCS feasible set can solve an optimization problem while producing an unsafe path; \cref{tab:tro_audited_gcs,tab:tro_soundness_audit} make this distinction explicit.

\Cref{tab:tro_experiment_matrix} lists the thirteen experiment artifacts used in the rewrite. The table is included to make the experimental scope auditable: package microbenchmarks, cache reuse, end-to-end planning, grow-stop ablation, GCS composition, merger safety, cross-robot transfer, a same-oracle single-query baseline, dynamic updates, parallel scaling, mechanism diagnostics, and paper-wide soundness accounting are generated by separate scripts and then consumed directly by the manuscript tables.

\begin{table}[t]
\centering
\caption{Experiment matrix for the TRO rewrite. Each artifact is generated by the listed script and consumed directly by the paper tables.}
\label{tab:tro_experiment_matrix}
\footnotesize
\resizebox{\columnwidth}{!}{%
\begin{tabular}{llll}
\toprule
ID & Script & Primary metric & Paper result \\
\midrule
Exp.1 & paper\_01\_epiaabb\_pipeline.py & endpoint iAABB source cost/tightness & tab:tro\_endpoint\_sources \\
Exp.2 & paper\_02\_link\_envelope\_pipeline.py & link envelope representation trade-off & tab:tro\_link\_envelopes \\
Exp.3 & paper\_03\_marcucci\_envelope\_build.py & LECT reuse and FFB depth sensitivity & tab:tro\_lect\_reuse, tab:tro\_ffb\_depth\_sweep \\
Exp.4 & paper\_04\_marcucci\_combined.py; paper\_04\_baselines\_marcucci.py; paper\_04\_rrt\_connect\_baseline.py; paper\_14\_shelf\_anytime\_tradeoff.py & Marcucci envelope variants and current live baseline trade-off & fig:tro\_shelf\_anytime\_tradeoff, tab:tro\_main\_shelf\_best\_tradeoff \\
Exp.5 & paper\_04\_grower\_tradeoff.py & grow-stop build/path-quality curve & tab:tro\_grower\_tradeoff \\
Exp.6 & paper\_04\_audited\_corridor\_gcs.py & audited corridor GCS audit pass rate & tab:tro\_audited\_gcs \\
Exp.7 & paper\_07\_merger\_protected\_study.py & direct vs protected merger/GCS safety & tab:tro\_merger\_control \\
Exp.8 & paper\_05\_random\_robot\_scenes.py & cross-robot random-scene SBF transfer & tab:tro\_random\_scenes \\
Exp.9 & paper\_06\_obstacle\_rebuild.py & dynamic obstacle deletion/rebuild & tab:tro\_dynamic\_rebuild \\
Exp.10 & paper\_07\_parallel\_scaling.py & thread sweep and scaling boundary & tab:tro\_parallel\_scaling \\
Exp.11 & paper\_11\_soundness\_audit\_suite.py & paper-wide strict audit accounting & tab:tro\_soundness\_audit \\
Exp.12 & paper\_12\_random\_scene\_iris\_np\_gcs\_baseline.py; paper\_12\_random\_scene\_rrt\_baseline.py; paper\_12\_random\_scene\_ompl\_baselines.py; paper\_15\_random\_anytime\_tradeoff.py & current random-scene baseline trade-off & fig:tro\_random\_anytime\_tradeoff, tab:tro\_main\_random\_best\_tradeoff \\
Exp.13 & paper\_13\_mechanism\_diagnostics.py & claim-boundary mechanism diagnostics & tab:tro\_mechanism\_diagnostics \\
\bottomrule
\end{tabular}
}
\end{table}


\subsection{Package-Level Envelope Evidence}\label{sec:exp_envelopes}

\Cref{tab:tro_endpoint_sources} isolates endpoint interval-AABB sources on the IIWA14 model. IFK gives inexpensive certified bounds but is relatively loose, whereas CritSample and Analytical bounds are tighter at higher cost. Monte Carlo sampling is reported as a reference-quality probe rather than a certificate. The comparison establishes why the paper-facing SBF configuration uses CritSample endpoint boxes for fast coverage experiments while retaining the strict-certificate mode as a conservative baseline.

\begin{table}[t]
\centering
\caption{Endpoint interval-envelope source microbenchmark averaged over the sampled width bins.}
\label{tab:tro_endpoint_sources}
\footnotesize
\resizebox{\columnwidth}{!}{%
\begin{tabular}{lrrrr}
\toprule
Source & Time (us) & Volume & Mean gap & Certified \\
\midrule
IFK & 1.437 & 0.5131 & 0 & yes \\
CritSample & 9.865 & 0.1429 & 0.0016 & no \\
Analytical & 1209.8 & 0.1448 & 1.87e-06 & yes \\
MC & 5234.1 & 0.1190 & 0.0212 & no \\
\bottomrule
\end{tabular}
}
\end{table}


\Cref{tab:tro_link_envelopes} then evaluates link-envelope representations using the same fixed box table. Increasing LinkIAABB subdivisions tightens the envelope with small microkernel overhead, while HullGrid variants trade additional payload and evaluation cost for tighter volumetric coverage. These measurements justify the two operating regimes used later: LinkIAABB for the fast paper-facing SBF pipeline, and HullGrid/LECT when cached geometric payloads are the scientific object under study.

\begin{table}[t]
\centering
\caption{Link-envelope representation trade-off. Ratios are relative to Crit+LinkIAABB S=1.}
\label{tab:tro_link_envelopes}
\footnotesize
\resizebox{\columnwidth}{!}{%
\begin{tabular}{lrrrr}
\toprule
Envelope & Eval (us) & Vol. ratio & Payload & D32 time (h) \\
\midrule
Crit+LinkIAABB S=1 & 12.35 & 1.000 & 48.0 B & 34.31 \\
Crit+LinkIAABB S=2 & 12.41 & 0.9045 & 48.0 B & 34.47 \\
Crit+LinkIAABB S=4 & 11.76 & 0.8625 & 48.0 B & 32.67 \\
Crit+LinkIAABB S=8 & 11.98 & 0.8429 & 48.0 B & 33.28 \\
Crit+HullGrid d=0.02 strict-pad & 166.83 & 0.6035 & 2.09 KB & 463.40 \\
Crit+HullGrid d=0.04 strict-pad & 66.12 & 0.7446 & 409.8 B & 183.68 \\
Crit+HullGrid d=0.06 strict-pad & 39.37 & 0.9073 & 187.6 B & 109.36 \\
\bottomrule
\end{tabular}
}
\end{table}


\subsection{LECT Reuse Across Build Protocols}\label{sec:exp_lect}

\Cref{tab:tro_lect_reuse} reports cold, warm, and cross-scene protocols for the envelope cache under the optimized full-run configuration. The diagnostic counters show that warm cache replay removes most endpoint-scoring work, but the wall-clock rows also show the important systems caveat: for cheap LinkIAABB payloads, validation and materialization overhead can dominate, so warm replay is not automatically faster than a fresh cold build. The key interpretation is therefore not that cached boxes alone solve the planning problem or always reduce every timing row; rather, LECT decouples reusable kinematic evidence from scene-specific collision filtering so that expensive envelope payloads can be amortized over repeated planning sessions.

\begin{table}[t]
\centering
\caption{LECT cache/build reuse under cold, warm, and cross-scene protocols.}
\label{tab:tro_lect_reuse}
\footnotesize
\resizebox{\columnwidth}{!}{%
\begin{tabular}{llrrrr}
\toprule
Variant & Protocol & Build (s) & Cache (MB) & Audit SR & Boxes \\
\midrule
coverage\_hybrid & cold & 0.7347 & 18.14 & 0.8000 & 129.70 \\
coverage\_hybrid & warm & 0.8455 & 21.26 & 0.8000 & 77.60 \\
coverage\_hybrid & cross\_scene & 0.9323 & 29.89 & 0.8000 & 67.60 \\
crit\_link\_coverage & cold & 0.2239 & 0.4896 & 0.8000 & 139.00 \\
crit\_link\_coverage & warm & 0.2944 & 0.5943 & 0.8000 & 138.10 \\
crit\_link\_coverage & cross\_scene & 0.3228 & 0.9324 & 0.8000 & 145.40 \\
\bottomrule
\end{tabular}
}
\end{table}


\subsection{End-to-End Shelf+IIWA Planning}\label{sec:exp_marcucci}

\Cref{tab:tro_marcucci_e2e} gives the main shelf+IIWA result. The current quality-aware configuration builds a connected SBF with mean build time \TroSbfBuildMean\,s and mean box count \TroSbfBoxMean, then answers the five canonical queries with audit success rate one. The path lengths reflect the conservative box-corridor representation: some routes are shorter than the v6 reference after local audited bridging, while others remain longer because the planner refuses to count un-audited shortcuts.

\begin{table}[t]
\centering
\caption{Shelf+IIWA end-to-end SBF result. The build row reports build time, mean boxes, and segment edges.}
\label{tab:tro_marcucci_e2e}
\footnotesize
\resizebox{\columnwidth}{!}{%
\begin{tabular}{lrrrrr}
\toprule
Item & SR & Audit SR & Time (s) & Length/Boxes & Repairs/Edges \\
\midrule
AS->TS & 1.000 & 1.000 & 1.75e-04 & 2.845 & 0 \\
TS->CS & 1.000 & 1.000 & 0.1016 & 2.338 & 1.000 \\
CS->LB & 1.000 & 1.000 & 0.5625 & 3.140 & 0 \\
LB->RB & 1.000 & 1.000 & 1.27e-04 & 4.429 & 0 \\
RB->AS & 1.000 & 1.000 & 2.03e-04 & 2.603 & 0 \\
Build summary & -- & -- & 0.4017 & 155.60 & 5.000 \\
\bottomrule
\end{tabular}
}
\end{table}


The grow-stop policy is evaluated in \cref{tab:tro_grower_tradeoff} and visualized in \cref{fig:tro_grower_tradeoff}. A purely time-minimal grow stop produces only about thirty connected boxes and behaves close to repaired BiRRT, which weakens both path quality and the algorithmic contribution. The numerical balance score reaches its fastest knee at a smaller floor, but the paper-facing default deliberately uses a denser 128-box floor: it preserves a visible reusable corridor while keeping the build below half a second. Higher floors add boxes and build time without removing the remaining topology-driven local repair events.

\begin{table}[t]
\centering
\caption{Grow-stop trade-off on the Marcucci scene without an imposed 0.5 s build cap. The selected floor is used by subsequent paper-facing SBF runs.}
\label{tab:tro_grower_tradeoff}
\footnotesize
\resizebox{\columnwidth}{!}{%
\begin{tabular}{rrrrrr}
\toprule
Quality floor & Build (s) & Boxes & Path sum & Repairs & Audit SR \\
\midrule
0 & 0.2382 & 29.90 & 15.27 & 1.000 & 1.000 \\
64 & 0.3099 & 87.20 & 15.00 & 1.000 & 1.000 \\
128 (selected) & 0.4492 & 155.00 & 15.42 & 1.000 & 1.000 \\
320 & 0.7419 & 349.70 & 15.58 & 1.000 & 1.000 \\
512 & 1.148 & 537.40 & 15.39 & 1.000 & 1.000 \\
\bottomrule
\end{tabular}
}
\end{table}


\begin{figure}[t]
  \centering
  \includegraphics[width=0.98\columnwidth]{figures/marcucci_grower_tradeoff.png}
  \caption{Grow-stop trade-off on the shelf+IIWA scene. The curve shows why stopping immediately after connectivity is too early, while very high quality floors add build time with limited path-quality gain.}
  \label{fig:tro_grower_tradeoff}
\end{figure}

\subsection{Audited GCS Composition}\label{sec:exp_gcs}

\Cref{tab:tro_merger_control} and \cref{tab:tro_audited_gcs} address the GCS question directly. The direct merger+GCS control solves all five shortest-path optimization instances but passes strict audit on \TroDirectGcsAuditPass\ of them. The failure mode is structural: the direct graph gives GCS access to provisional box interiors that were never certified as convex subsets of $\Cfree$. A second issue is purely geometric: LinearGCS requires adjacent convex regions to have a feasible overlap, whereas SBF segment-edges may connect two conservative boxes by an audited segment even when the boxes themselves do not intersect. We therefore add an overlap-expanded SBF-corridor adapter that expands consecutive query boxes just enough to create pairwise GCS overlap; this adapter passes strict audit on \TroExpandedGcsAuditPass\ solved expanded attempts and is still not counted unless the final path is audited. The final protected interface either uses that audited expanded corridor or a narrow audited path-tube corridor, and under this interface audited GCS passes \TroAuditedGcsAuditPass\ on the five Marcucci queries.

\begin{table}[t]
\centering
\caption{Merger and protected-GCS study. Merging or expanding boxes changes GCS feasibility, not the safety certificate.}
\label{tab:tro_merger_control}
\footnotesize
\resizebox{\columnwidth}{!}{%
\begin{tabular}{llrrrr}
\toprule
Mode & Regions & Solve & Audit & Unsafe & Max gap \\
\midrule
direct\_merger\_gcs & merged/provisional SBF boxes & 5 & 0 & 5 & -- \\
sbf\_box\_corridor & raw SBF query boxes & 2 & 0 & 2 & -- \\
overlap\_expanded\_sbf\_corridor & pairwise overlap-expanded query boxes & 4 & 0 & 4 & 0.7853 \\
audited\_path\_tube & narrow audited path-tube boxes & 5 & 5 & 0 & -- \\
protected\_final\_policy & first strict-audited GCS attempt, else audited SBF fallback & 11 & 5 & 0 & -- \\
\bottomrule
\end{tabular}
}
\end{table}

\begin{table}[t]
\centering
\caption{GCS composition result. Expanded boxes repair LinearGCS overlap feasibility but are counted only after strict audit.}
\label{tab:tro_audited_gcs}
\footnotesize
\resizebox{\columnwidth}{!}{%
\begin{tabular}{lrrrrrrr}
\toprule
Query & Direct solved & Direct audit & Expanded solved & Expanded audit & Final audit & Final source & Length \\
\midrule
AS->TS & yes & no & yes & no & yes & gcs\_strict\_audited & 2.800 \\
TS->CS & yes & no & no & no & yes & gcs\_strict\_audited & 2.313 \\
CS->LB & yes & no & yes & no & yes & gcs\_strict\_audited & 3.121 \\
LB->RB & yes & no & yes & no & yes & gcs\_strict\_audited & 4.409 \\
RB->AS & yes & no & yes & no & yes & gcs\_strict\_audited & 2.543 \\
\bottomrule
\end{tabular}
}
\end{table}


This result should be read as a composition boundary rather than as a replacement for IRIS-style certified convex decomposition. SBF boxes are useful for quickly finding reusable corridors. If those corridors are to be consumed by a convex optimizer, non-overlapping segment-edge boxes must first be made GCS-feasible by a protected overlap adapter or by a path-tube construction, and the optimizer output must still be audited. Expanding a box boundary repairs LinearGCS feasibility; it does not by itself create a free-space certificate.

\subsection{Generalization, Updates, and Scaling}\label{sec:exp_generalization}

\Cref{tab:tro_random_scenes} evaluates the standalone SBF implementation on randomized UR5 and Panda scenes under a bounded 3\,s per-case build budget. The purpose is narrower than the shelf+IIWA study: it checks that the envelope/build/query stack is not specialized to one robot or one obstacle layout, while leaving difficult random cases as benchmark failures rather than hiding them behind unbounded build time. Because the standalone migration does not include Drake or OMPL baselines, \cref{tab:tro_random_rrt_baseline} adds a same-oracle RRT-Connect comparison~\cite{kuffner2000rrt} on the identical random-scene protocol. This baseline has no reusable build phase and should be read as a single-query reference, not as a multi-query cache competitor.

\begin{table}[t]
\centering
\caption{Cross-robot random-scene standalone SBF evaluation.}
\label{tab:tro_random_scenes}
\footnotesize
\resizebox{\columnwidth}{!}{%
\begin{tabular}{lllrrrr}
\toprule
Robot & Difficulty & Method & Build (s) & Boxes & Query SR & Audit SR \\
\midrule
ur5 & easy & SBF-SH & 0.4742 & 310.20 & 1.000 & 1.000 \\
ur5 & medium & SBF-SH & 0.5279 & 208.40 & 1.000 & 1.000 \\
panda & easy & SBF-SH & 1.987 & 518.00 & 1.000 & 1.000 \\
panda & medium & SBF-SH & 2.058 & 513.60 & 1.000 & 1.000 \\
\bottomrule
\end{tabular}
}
\end{table}

\begin{table}[t]
\centering
\caption{Raw OMPL RRTConnect single-query baseline on the random-scene protocol. It uses the SBF collision checker, no reusable build, and no shortcut simplification by default.}
\label{tab:tro_random_rrt_baseline}
\footnotesize
\resizebox{\columnwidth}{!}{%
\begin{tabular}{lllrrrrr}
\toprule
Robot & Difficulty & Planner & Build (s) & Query (s) & Audit SR & Path & Boxes \\
\midrule
panda & easy & SBF-SH & 1.987 & 0.0022 & 1.000 & 5.432 & 518.00 \\
panda & easy & OMPL RRTConnect & 0 & 0.0017 & 1.000 & 4.871 & -- \\
panda & medium & SBF-SH & 2.058 & 0.0028 & 1.000 & 4.821 & 513.60 \\
panda & medium & OMPL RRTConnect & 0 & 0.0019 & 1.000 & 5.718 & -- \\
ur5 & easy & SBF-SH & 0.4742 & 0.0019 & 1.000 & 5.493 & 310.20 \\
ur5 & easy & OMPL RRTConnect & 0 & 0.0017 & 1.000 & 5.513 & -- \\
ur5 & medium & SBF-SH & 0.5279 & 0.2156 & 1.000 & 7.348 & 208.40 \\
ur5 & medium & OMPL RRTConnect & 0 & 0.0026 & 1.000 & 6.900 & -- \\
\bottomrule
\end{tabular}
}
\end{table}


\Cref{tab:tro_dynamic_rebuild} studies obstacle insertion after an SBF build. The reported rebuild is localized: boxes colliding with the inserted obstacle are deleted and adjacency is recomputed. This is not a full coverage regrowth, so it measures the cost of maintaining a conservative box graph under a local scene edit.

\begin{table}[t]
\centering
\caption{Dynamic obstacle insertion and localized box deletion/rebuild cost with post-update query/audit accounting when available.}
\label{tab:tro_dynamic_rebuild}
\footnotesize
\resizebox{\columnwidth}{!}{%
\begin{tabular}{lrrrrrrr}
\toprule
Obstacle & Initial build (s) & Boxes & Removed & Removal & Rebuild (s) & Query SR & Audit SR \\
\midrule
front\_bin & 0.1087 & 64.00 & 56.00 & 0.8750 & 1.54e-04 & -- & -- \\
\bottomrule
\end{tabular}
}
\end{table}


\Cref{tab:tro_parallel_scaling} reports a thread sweep for the end-to-end Marcucci pipeline. The current implementation obtains reliable audit success at every thread count, but the small problem size and synchronization overhead prevent monotone speedup. This negative scaling result is included because it bounds the present engineering claim: the method is fast in absolute build time on the benchmark, but the current standalone grower is not yet a strong-scaling implementation.

\begin{table}[t]
\centering
\caption{Parallel build scaling for the Marcucci end-to-end pipeline.}
\label{tab:tro_parallel_scaling}
\footnotesize
\resizebox{\columnwidth}{!}{%
\begin{tabular}{rrrrrr}
\toprule
Threads & Build (s) & Speedup & Efficiency & Boxes & Audit SR \\
\midrule
1 & 0.4145 & 1.000 & 1.000 & 137.60 & 1.000 \\
2 & 0.5866 & 0.7066 & 0.3533 & 149.40 & 1.000 \\
4 & 0.3958 & 1.047 & 0.2618 & 149.80 & 1.000 \\
8 & 0.3944 & 1.051 & 0.1314 & 155.20 & 1.000 \\
16 & 0.4067 & 1.019 & 0.0637 & 155.20 & 1.000 \\
\bottomrule
\end{tabular}
}
\end{table}


\Cref{tab:tro_mechanism_diagnostics} collects the claim-boundary mechanisms that should not be hidden behind a single success-rate number. It reports the same-oracle random-scene baseline range, the grow-stop knee, the GCS negative control, the parallel-scaling limit, and the paper-wide audit count. These rows are not additional claims; they are guardrails on how the main results should be interpreted.

\begin{table}[t]
\centering
\caption{Mechanism diagnostics for claim boundaries, failure modes, and negative controls.}
\label{tab:tro_mechanism_diagnostics}
\footnotesize
\resizebox{\columnwidth}{!}{%
\begin{tabular}{llll}
\toprule
Scope & Metric & Value & Interpretation \\
\midrule
shelf baseline & OMPL RRTConnect audit SR / median time & 1.00; 0.010 s & OMPL RRTConnect is a single-query shelf reference with zero reusable build cost \\
random-scene baseline & OMPL RRTConnect audit SR range & 1.00-1.00 & random-scene SBF rows are now contextualized against an OMPL single-query planner using the same SBF collision checker \\
grow-stop policy & quality floor 0 vs selected knee & boxes 29.9 to 87.2; total path 15.269 to 14.996 & larger connected-box floor improves reusable corridor quality but does not remove all local repairs on this topology \\
GCS composition & direct / expanded / protected audit & 0/5; 0/4; 5/5 & box expansion repairs overlap feasibility but safety is recovered only by the protected audited interface \\
parallel scaling & best speedup / lowest efficiency & 1.00x at 1 threads; 0.12 efficiency at 8 threads & small Marcucci builds are dominated by serial scheduling, bridge/refine, and merge overhead; parallelism is not claimed as strong scaling \\
paper-wide audit & strict audit pass & 663/668 & failed random-scene queries remain failures; solved-but-unsafe paths are separated as negative controls \\
\bottomrule
\end{tabular}
}
\end{table}


\subsection{Soundness Accounting}\label{sec:exp_soundness}

\Cref{tab:tro_soundness_audit} summarizes the safety accounting across the paper-facing path sets. The table separates raw solver success, strict audit success, local repair or fallback use, solved-but-unsafe cases, and remaining unsafe assumptions. Across all aggregated full-run paper path artifacts, \TroPaperAuditPass\ paths pass the strict audit accounting used in the tables, and the solved-but-unsafe GCS attempts are retained only as negative controls. This convention is stricter than reporting a single planner success rate, but it is necessary for a method whose central claim is conservative reuse of geometric evidence.

\begin{table}[t]
\centering
\caption{Paper-wide soundness audit suite. Solved-but-unsafe paths are counted separately and never as planning success.}
\label{tab:tro_soundness_audit}
\footnotesize
\resizebox{\columnwidth}{!}{%
\begin{tabular}{lrrrrl}
\toprule
Scope & Paths & Audit SR & Repair/Fallback & Solved unsafe & Residual assumption \\
\midrule
Exp.4 Marcucci SBF & 50 & 1.000 & 10.00 & 0 & none after strict audit \\
Exp.3 grow-stop sweep & 450 & 1.000 & 90.00 & 0 & none after strict audit \\
Exp.7 audited corridor GCS & 5 & 1.000 & 2.000 & 8 & none for final counted paths; 8 unsafe solved attempts rejected \\
Exp.7 direct merger+GCS control & 5 & 0 & 0 & 5 & provisional/merged region interiors are not convex-free certificates \\
Exp.4 OMPL RRTConnect shelf baseline & 50 & 1.000 & 0 & 0 & failed baseline trials remain failures \\
Exp.5 random robot scenes & 8 & 1.000 & 5.000 & 0 & none for audited successes; failures remain failures \\
Exp.5 OMPL RRTConnect random-scene baseline & 40 & 1.000 & 0 & 0 & failed baseline trials remain failures \\
Exp.8 parallel sweep & 60 & 1.000 & 12.00 & 0 & none for audited successes \\
\bottomrule
\end{tabular}
}
\end{table}

